% Created 2026-06-22 Mon 13:07
% Intended LaTeX compiler: xelatex
\documentclass[11pt]{article}
\usepackage{graphicx}
\usepackage{longtable}
\usepackage{wrapfig}
\usepackage{rotating}
\usepackage[normalem]{ulem}
\usepackage{capt-of}
\usepackage{hyperref}
\author{Bazinevis}
\date{\today}
\title{}
\hypersetup{
 pdfauthor={Bazinevis},
 pdftitle={},
 pdfkeywords={},
 pdfsubject={},
 pdfcreator={Emacs 30.2 (Org mode 9.7.11)}, 
 pdflang={English}}
\begin{document}

\tableofcontents

\section{Chapter 1: Axioms of Probability}
\label{sec:org754a5a2}
\subsection{Axiom of Choice (AC)}
\label{sec:org8b6f3b7}
\subsubsection{Informal Idea}
\label{sec:org2b36185}
The Axiom of Choice states that:

\begin{quote}
Given any collection of nonempty sets, it is possible to choose one element from each set, even if no explicit rule for making the choices is known.
\end{quote}

In simple terms:
\begin{itemize}
\item If you have finitely many sets, choosing one element from each is easy.
\item If you have infinitely many sets, especially infinitely many arbitrary sets, the existence of a choice function is not obvious.
\item The Axiom of Choice guarantees that such a function exists.
\end{itemize}
\subsubsection{Formal Statement}
\label{sec:orgf49d436}
Let

\[
\mathcal{F} = \{A_i : i \in I\}
\]

be a collection of nonempty sets. Then there exists a function

\[
f : I \to \bigcup_{i\in I} A_i
\]

such that

\[
f(i) \in A_i
\]

for every \(i \in I\). The function \(f\) is called a \textbf{choice function}.
\subsubsection{Example: Finite Collection}
\label{sec:orgc98abb5}
Suppose

\[
A_1=\{1,2\}
\]
\[
A_2=\{a,b,c\}
\]
\[
A_3=\{x,y\}
\]

A choice function could be:

\[
f(1)=1,\quad
f(2)=b,\quad
f(3)=x
\]

Here we selected one element from each set. No axiom is really needed because there are only finitely many choices.
\subsubsection{Example: Infinite Collection}
\label{sec:org36b84bb}
Consider the family

\[
A_n=\{n,n+1\}
\]

for every natural number \(n\). A choice function might be

\[
f(n)=n.
\]

Again, a clear rule exists. The interesting cases occur when no obvious rule exists.
\subsubsection{Why Is It Controversial?}
\label{sec:orge7f0a48}

Imagine an infinite collection of nonempty sets where:

\begin{itemize}
\item The sets are arbitrary.
\item No formula distinguishes one element from another.
\item No explicit selection procedure is available.
\end{itemize}

The Axiom of Choice says:

\begin{quote}
A choice function still exists.
\end{quote}

Even if nobody can describe it.

Many mathematicians initially found this surprising because it asserts existence without providing a construction.
\subsubsection{Equivalent Statements}
\label{sec:orgaf8d0de}

The following famous results are equivalent to the Axiom of Choice.
\begin{enumerate}
\item Well-Ordering Theorem
\label{sec:org3aa999d}

Every set can be well-ordered.

That means:

Every set can be arranged so that every nonempty subset has a smallest element.

Even the real numbers can be well-ordered, although no explicit well-ordering is known.
\item Zorn's Lemma
\label{sec:orga2bfb95}

If every chain in a partially ordered set has an upper bound, then the set contains a maximal element.

Zorn's Lemma is heavily used in algebra.

Examples:

\begin{itemize}
\item Every vector space has a basis.
\item Every ring has a maximal ideal.
\end{itemize}
\item Tychonoff's Theorem
\label{sec:orgc44967e}

The product of any collection of compact topological spaces is compact.

This is one of the most important theorems in topology.
\end{enumerate}
\subsubsection{Consequences}
\label{sec:org80d643b}

Using the Axiom of Choice, mathematicians can prove:

\begin{itemize}
\item Every vector space has a basis.
\item Every set can be well-ordered.
\item Infinite products of nonempty sets are nonempty.
\item Many existence theorems in algebra and analysis.
\end{itemize}
\subsubsection{The Vitali Set}
\label{sec:org136ab4f}

One of the most famous consequences is the existence of a Vitali set.

A Vitali set is constructed by selecting exactly one representative from each equivalence class of the relation

\[
x \sim y
\]

if and only if

\[
x-y \in \mathbb{Q}.
\]

There are uncountably many equivalence classes.

To choose one representative from each class, we need the Axiom of Choice.

The resulting set is not Lebesgue measurable.

Thus:

\begin{itemize}
\item Axiom of Choice
-> Vitali set exists.
\item Vitali set
-> Non-measurable sets exist.
\end{itemize}
\subsubsection{Why Measure Theory Cares}
\label{sec:org197e0c6}

Without the Axiom of Choice:

\begin{itemize}
\item It is possible that every subset of the real numbers is measurable.
\end{itemize}

With the Axiom of Choice:

\begin{itemize}
\item Non-measurable sets exist.
\item The Vitali set is the standard example.
\end{itemize}

This is why the Axiom of Choice appears frequently in advanced measure theory.
\subsubsection{Intuition}
\label{sec:org5105cbc}

The Axiom of Choice does \textbf{not} tell us how to make the choices.

It only says that a choice function exists.

Think of having infinitely many boxes, each containing at least one object.

The Axiom of Choice says:

\begin{quote}
You can pick one object from every box simultaneously,
even if no rule tells you which object to pick.
\end{quote}
\subsubsection{Summary}
\label{sec:org82e85bd}

\begin{itemize}
\item The Axiom of Choice asserts the existence of a choice function.
\item It allows choosing one element from each set in an arbitrary collection of
nonempty sets.
\item It is equivalent to:
\begin{itemize}
\item Well-Ordering Theorem
\item Zorn's Lemma
\item Tychonoff's Theorem
\end{itemize}
\item It implies the existence of non-measurable sets such as the Vitali set.
\item Many fundamental results in modern mathematics depend on it.
\end{itemize}
\subsection{Measure Theory}
\label{sec:orgbfa3f46}
\subsubsection{Motivation}
\label{sec:org1e1d359}

Measure theory is the mathematical study of:

\begin{itemize}
\item Length
\item Area
\item Volume
\item Probability
\end{itemize}

and, more generally, the size of sets.

Its goal is to answer questions such as:

\begin{itemize}
\item How long is a subset of the real line?
\item How large is a subset of the plane?
\item What is the probability of an event?
\item Can every subset of the real numbers be assigned a size?
\end{itemize}

Measure theory provides a rigorous framework for answering these questions.
\subsubsection{Historical Background}
\label{sec:org9a4394a}

Classical geometry measures simple objects:

\begin{itemize}
\item Intervals
\item Rectangles
\item Circles
\item Cubes
\end{itemize}

However, mathematicians discovered sets that are much more complicated.

Examples:

\begin{itemize}
\item Fractals
\item Countable dense sets
\item Non-measurable sets
\end{itemize}

A more general theory was needed.

Measure theory was developed primarily by \textbf{Henri Lebesgue} (1875–1941), whose ideas generalized the notions of length, area, and integration and laid the foundation for modern analysis and probability theory.
\subsubsection{Basic Idea}
\label{sec:org2d76469}

A measure assigns a nonnegative number to a set.

Think of a measure as a generalized notion of size.

For a set A:

\[
\mu(A)
\]

means

\begin{quote}
"The size of A."
\end{quote}

Depending on the context, size may mean:

\begin{itemize}
\item Length
\item Area
\item Volume
\item Probability
\item Mass
\end{itemize}
\subsubsection{Example: Length on the Real Line}
\label{sec:orged313b9}

Consider the interval

\[
[0,5].
\]

Its length is

\[
\mu([0,5]) = 5.
\]

Similarly,

\[
\mu([2,7]) = 5.
\]

Measure theory generalizes this notion to much more complicated sets.
\subsubsection{Why We Need Sigma-Algebras}
\label{sec:org2adc731}

Not every collection of sets behaves nicely.

We want a collection of sets that is closed under operations such as:

\begin{itemize}
\item Complements
\item Countable unions
\item Countable intersections
\end{itemize}

Such a collection is called a sigma-algebra.
\subsubsection{Definition: Sigma-Algebra}
\label{sec:org48d4519}

Let X be a set.

A collection

\[
\mathcal F
\]

of subsets of X is a sigma-algebra if:
\begin{enumerate}
\item Contains the whole space
\label{sec:org0235403}

\[
X \in \mathcal F.
\]
\item Closed under complements
\label{sec:orgf4ca345}

If

\[
A \in \mathcal F
\]

then

\[
A^c \in \mathcal F.
\]
\item Closed under countable unions
\label{sec:org154c357}

If

\[
A_1,A_2,A_3,\ldots \in \mathcal F
\]

then

\[
\bigcup_{n=1}^{\infty} A_n
\in \mathcal F.
\]
\end{enumerate}
\subsubsection{Why Countable Operations?}
\label{sec:orge11dc2f}

Countable operations are sufficient for:

\begin{itemize}
\item Infinite sums
\item Limits
\item Convergence
\item Probability theory
\end{itemize}

Allowing arbitrary unions creates technical problems.

Measure theory therefore focuses on countable operations.
\subsubsection{Measurable Space}
\label{sec:orgd5e7814}

A measurable space consists of:

\[
(X,\mathcal F).
\]

where

\begin{itemize}
\item X is the set of possible outcomes.
\item \(\mathcal F\) is a sigma-algebra.
\end{itemize}

The sets in \(\mathcal F\) are called measurable sets.

Only measurable sets are assigned measures.
\subsubsection{Measure}
\label{sec:org82f7442}

A measure is a function

\[
\mu : \mathcal F \to [0,\infty].
\]

satisfying:
\begin{enumerate}
\item Nonnegativity
\label{sec:orge84a250}

\[
\mu(A)\ge 0.
\]
\item Empty Set
\label{sec:org9d7f7dc}

\[
\mu(\emptyset)=0.
\]
\item Countable Additivity
\label{sec:orgbf37dd0}

If

\[
A_1,A_2,A_3,\ldots
\]

are pairwise disjoint, then

\[
\mu\left(
\bigcup_{n=1}^{\infty} A_n
\right)
=
\sum_{n=1}^{\infty}
\mu(A_n).
\]

This is the most important property.
\end{enumerate}
\subsubsection{Measure Space}
\label{sec:org1d7055f}

A measure space consists of:

\[
(X,\mathcal F,\mu).
\]

where:

\begin{itemize}
\item X is a set.
\item \(\mathcal F\) is a sigma-algebra.
\item \(\mu\) is a measure.
\end{itemize}

This is the fundamental object of measure theory.
\subsubsection{Example: Counting Measure}
\label{sec:orgad31b13}

Let

\[
X=\mathbb N.
\]

Define

\[
\mu(A)
=
\text{number of elements of }A.
\]

Examples:

\[
\mu(\{1,5,8\})=3.
\]

\[
\mu(\mathbb N)=\infty.
\]

This is called the counting measure.
\subsubsection{Example: Probability Measure}
\label{sec:org3e1f5ae}

A probability measure satisfies

\[
P(X)=1.
\]

Probability theory is therefore a special case of measure theory.

Correspondence:

\begin{center}
\begin{tabular}{ll}
Measure Theory & Probability\\
\hline
X & Sample Space\\
Measurable Set & Event\\
Measure \(\mu\) & Probability P\\
\(\mu(X)\) & 1\\
\end{tabular}
\end{center}
\subsubsection{Lebesgue Measure}
\label{sec:orge36493e}

The most important measure on the real line is the Lebesgue measure.

For ordinary intervals:

\[
m([a,b])=b-a.
\]

Examples:

\[
m([0,1])=1
\]

\[
m([2,7])=5
\]

Lebesgue measure extends the notion of length to very complicated sets.
\subsubsection{Null Sets}
\label{sec:orgee51c95}

A set with measure zero is called a null set.

Examples:
\begin{enumerate}
\item A Single Point
\label{sec:org910327a}

\[
m(\{5\})=0.
\]
\item Finite Sets
\label{sec:orgdf20a4d}

\[
m(\{1,2,3\})=0.
\]
\item Countable Sets
\label{sec:orge0d7468}

\[
m(\mathbb Q)=0.
\]

Even though the rational numbers are dense in the real line, their total
Lebesgue measure is zero.
\end{enumerate}
\subsubsection{Why Rational Numbers Have Measure Zero}
\label{sec:orge6d463a}

Each rational number has length zero.

Since the rationals are countable, we can cover them by intervals whose total
length is arbitrarily small.

Therefore

\[
m(\mathbb Q)=0.
\]

This is one of the first surprising results in measure theory.
\subsubsection{Non-Measurable Sets}
\label{sec:org6fc22ba}

One might hope every subset of the real numbers has a measure.

Unfortunately this is false.

Using the Axiom of Choice, one can construct:

\begin{itemize}
\item Vitali sets
\item Other non-measurable sets
\end{itemize}

These sets cannot consistently be assigned a Lebesgue measure.
\subsubsection{Lebesgue Measurable Sets}
\label{sec:org2d87f70}

A set is Lebesgue measurable if it behaves well with respect to length.

The collection of all Lebesgue measurable sets forms a sigma-algebra.

This sigma-algebra contains:

\begin{itemize}
\item Open sets
\item Closed sets
\item Countable unions
\item Countable intersections
\item Borel sets
\item Many more complicated sets
\end{itemize}
\subsubsection{Integration}
\label{sec:org91528cb}

Measure theory leads naturally to integration.

Classical calculus uses the Riemann integral.

Lebesgue developed a more powerful integral.

Advantages:

\begin{itemize}
\item Handles more functions.
\item Works better with limits.
\item Dominates modern probability theory.
\item Essential in functional analysis.
\end{itemize}
\subsubsection{Connection to Probability}
\label{sec:org925b138}

Probability theory can be viewed as measure theory where

\[
\mu(X)=1.
\]

A probability space is

\[
(\Omega,\mathcal F,P).
\]

where:

\begin{itemize}
\item \(\Omega\) is the sample space.
\item \(\mathcal F\) is a sigma-algebra.
\item P is a probability measure.
\end{itemize}

Thus probability is essentially measure theory with total mass equal to one.
\subsubsection{Big Picture}
\label{sec:orga73799d}

Measure Theory asks:

\begin{quote}
How can we rigorously define the size of very complicated sets?
\end{quote}

To answer this, it introduces:

\begin{itemize}
\item Sigma-algebras
\item Measures
\item Measurable sets
\item Lebesgue measure
\item Integration
\end{itemize}

These ideas form the foundation of:

\begin{itemize}
\item Probability Theory
\item Functional Analysis
\item Ergodic Theory
\item Stochastic Processes
\item Modern Analysis
\end{itemize}
\subsubsection{Mental Map}
\label{sec:org2759fac}

Set X
  ->
Sigma Algebra \(\mathcal F\)
  ->
Measure \(\mu\)
  ->
Measure Space \((X,\mathcal F,\mu)\)
  ->
Lebesgue Measure
  ->
Lebesgue Integration
  ->
Probability Theory
  ->
Modern Analysis
\subsubsection{Summary}
\label{sec:org0ec55d8}

\begin{itemize}
\item Measure theory generalizes the idea of length, area, volume, and probability.
\item A sigma-algebra specifies which sets can be measured.
\item A measure assigns a size to measurable sets.
\item A measure space is \((X,\mathcal F,\mu)\).
\item Lebesgue measure generalizes ordinary length.
\item Countable sets such as the rationals have measure zero.
\item Some sets (e.g. Vitali sets) are not measurable.
\item Probability theory is a special case of measure theory.
\item Lebesgue integration is one of the central achievements of measure theory.
\end{itemize}
\subsection{Probability and Uncertainty}
\label{sec:org7794707}

\begin{itemize}
\item Probability is used when the outcome of an experiment is uncertain.
\item If we had complete information about all relevant factors, the outcome would be fully determined and probability would not be needed.
\item In practice, we rarely possess complete knowledge of a system.
\item Probability provides a mathematical framework for reasoning under uncertainty.
\end{itemize}

Examples:
\begin{itemize}
\item A coin toss follows physical laws, but we do not know the exact initial conditions.
\item Weather systems are governed by physics, but complete information is unavailable.
\item Therefore probability quantifies uncertainty arising from incomplete knowledge or practical unpredictability.
\end{itemize}
\subsection{1.1 Introduction}
\label{sec:orgf558ab6}

\subsubsection{Random, Certain, and Impossible Events}
\label{sec:org2249ee1}

\begin{itemize}
\item \textbf{Random Event}: May or may not occur.
\begin{itemize}
\item Example: A radioactive atom may or may not decay during a given time interval.
\item Example: An inspector may or may not find a defect in a product.
\end{itemize}

\item \textbf{Certain Event}: Guaranteed to occur.
\begin{itemize}
\item Example: Lightning is observed before thunder is heard.
\end{itemize}

\item \textbf{Impossible Event}: Can never occur.
\begin{itemize}
\item Example: An object traveling faster than the speed of light.
\end{itemize}
\end{itemize}
\subsubsection{Why Probability Is Needed}
\label{sec:org543cf23}

Knowing that an event is random is not enough. We want to measure \textbf{how likely} it is to occur.

Observation:
\begin{itemize}
\item When an experiment is repeated many times, the proportion of occurrences of an event often approaches a constant value.
\item Probability is used to describe this long-run behavior.
\end{itemize}

Example:
\begin{itemize}
\item Repeatedly tossing a fair coin results in the proportion of heads approaching 1/2.
\item Therefore:
\end{itemize}

\[
P(\text{Heads})=\frac{1}{2}
\]
\subsubsection{Historical Background}
\label{sec:org1388833}

Probability originated from the study of games of chance and gambling.

Important contributors:

\begin{itemize}
\item \href{../../_ocean_/people.org}{Luca Pacioli}
\item \href{../../_ocean_/people.org}{Niccolò Tartaglia}
\item \href{../../_ocean_/people.org}{Girolamo Cardano}
\item \href{../../_ocean_/people.org}{Galileo Galilei}
\end{itemize}

Major milestone:

\begin{itemize}
\item \href{../../_ocean_/people.org}{Blaise Pascal}
\item \href{../../_ocean_/people.org}{Pierre de Fermat}
\item \href{../../_ocean_/people.org}{Christiaan Huygens}
\end{itemize}

Further development:

\begin{itemize}
\item \href{../../_ocean_/people.org}{Jacob Bernoulli}
\item \href{../../_ocean_/people.org}{Abraham de Moivre}
\item \href{../../_ocean_/people.org}{Pierre-Simon Laplace}
\item \href{../../_ocean_/people.org}{Siméon Denis Poisson}
\item \href{../../_ocean_/people.org}{Carl Friedrich Gauss}
\item \href{../../_ocean_/people.org}{Pafnuty Chebyshev}
\item \href{../../_ocean_/people.org}{Andrey Markov}
\item \href{../../_ocean_/people.org}{Aleksandr Lyapunov}
\end{itemize}
\subsubsection{Interpretations of Probability}
\label{sec:orgca72181}

\begin{enumerate}
\item Classical Interpretation
\label{sec:org5e1de70}

Probability is defined as

\[
P(A)=
\frac{\text{number of favorable outcomes}}
     {\text{number of equally likely outcomes}}
\]

\begin{itemize}
\item Originated from gambling problems.
\item Applicable when outcomes are equally likely.
\end{itemize}

Example:

\[
P(\{6\})=\frac{1}{6}
\]

for a fair die.
\item Relative Frequency Interpretation
\label{sec:orgcb8ba7d}

Probability is the limiting relative frequency of occurrence in repeated trials.

\[
P(A)=
\lim_{n\to\infty}
\frac{n(A)}{n}
\]

where

\begin{itemize}
\item \(n\) = number of repetitions.
\item \(n(A)\) = number of times event \(A\) occurs.
\end{itemize}
\item Subjective Interpretation
\label{sec:org6ab1c8d}

\begin{itemize}
\item Probability represents a degree of belief based on available information.
\item Different individuals may assign different probabilities if they possess different information.
\end{itemize}

Examples:

\begin{itemize}
\item Probability a startup succeeds.
\item Probability a candidate wins an election.
\item Probability a patient responds to a treatment.
\end{itemize}
\end{enumerate}
\subsubsection{Relative Frequency Interpretation}
\label{sec:orge4fea1b}

\begin{quote}
\[
P(A)=
\lim_{n\to\infty}
\frac{n(A)}{n}
\]
\end{quote}
\subsubsection{Problems with the Frequency Definition}
\label{sec:org609c733}

\begin{enumerate}
\item Infinite experiments cannot be performed in practice.
\item The limit may not exist.
\item Subjective probabilities cannot be justified.
\end{enumerate}
\subsubsection{Hilbert's Challenge}
\label{sec:orgee98e36}

\href{../../_ocean_/people.org}{David Hilbert} proposed axiomatizing probability in 1900.
\subsubsection{Kolmogorov's Axiomatic Approach (1933)}
\label{sec:org8118f60}

\href{../../_ocean_/people.org}{Andrey Kolmogorov} built the modern axiomatic foundation of probability.
\subsubsection{Main Takeaway}
\label{sec:org79107fc}

Probability provides a rigorous framework for uncertainty using axioms.
\subsubsection{1.2 Sample Space and Events}
\label{sec:orga51b07c}

\begin{enumerate}
\item Sample Space
\label{sec:orgbebddbd}

\begin{itemize}
\item The \textbf{sample space} (\(S\)) is the set of all possible outcomes of an experiment.
\item Each outcome is called a \textbf{sample point}.
\end{itemize}

Examples:

\[
S=\{H,T\}
\]

\[
S=\{x:x\ge0\}
\]
\item Events
\label{sec:orgf98d4dd}

\begin{itemize}
\item An \textbf{event} is any \textbf{subset} of the sample space.
\item An event occurs when the observed outcome \textbf{belongs} to that subset.
\item Sample space is a set as well as event. If one member of a event set happens, we say that event happened. [hey ai add example here]
\end{itemize}

Examples:

\begin{itemize}
\item Coin toss: \(\{H\}\) = getting heads.
\item Light bulb lifetime ≥ 100 hours.
\item Bus arrives with 27 passengers between specific times.
\end{itemize}
\item Different Representations
\label{sec:org7aab8ae}

\begin{itemize}
\item The same experiment can have different sample-space representations.
\end{itemize}

Example:

\begin{itemize}
\item Bus arrival time measured in hours.
\item Bus arrival time measured in minutes after 11:00.
\end{itemize}

Different representations describe the same outcomes.
\item Basic Event Operations
\label{sec:org3c2b94e}

\begin{enumerate}
\item Subset
\label{sec:orgf92569e}

\[
E\subseteq F
\]

Whenever \(E\) occurs, \(F\) also occurs.
\item Equality
\label{sec:orgcf012ac}

\[
E=F
\]

iff

\[
E\subseteq F
\quad\text{and}\quad
F\subseteq E
\]
\item Intersection
\label{sec:orgd4f4270}

\[
E\cap F
\]

Occurs when both \(E\) and \(F\) occur.
\item Union
\label{sec:org89ed325}

\[
E\cup F
\]

Occurs when at least one of \(E\) or \(F\) occurs.
\item Complement
\label{sec:org446bb2f}

\[
E^c
\]

Occurs when \(E\) does not occur.
\item Difference
\label{sec:org7612965}

\[
E-F=E\cap F^c
\]

Occurs when \(E\) occurs but \(F\) does not.
\item Certain Event
\label{sec:org6ed6b7b}

\begin{itemize}
\item Event that always occurs.
\item The sample space \(S\).
\end{itemize}
\item Impossible Event
\label{sec:orgdcc9ca5}

\[
\varnothing
\]
\item Mutually Exclusive Events
\label{sec:org0a10df3}

\[
E\cap F=\varnothing
\]
\end{enumerate}
\item Operations on Collections of Events
\label{sec:org45e0120}

\begin{enumerate}
\item Union of Many Events
\label{sec:org836a6e1}

\[
\bigcup_i E_i
\]

At least one \(E_i\) occurs.
\item Intersection of Many Events
\label{sec:orgc070a8b}

\[
\bigcap_i E_i
\]

All \(E_i\) occur simultaneously.
\end{enumerate}
\item Example: Airport Queue
\label{sec:org2ecfcfe}

Let

\begin{itemize}
\item \(E\) = at least 5 planes waiting.
\item \(F\) = at most 3 planes waiting.
\item \(H\) = exactly 2 planes waiting.
\end{itemize}

Results:

\begin{itemize}
\item \(E^c\) = at most 4 planes waiting.
\item \(F^c\) = at least 4 planes waiting.
\item \(E\subseteq F^c\)
\item \(H\subseteq F\)
\item \(E\) and \(F\) are mutually exclusive.
\item \(E\) and \(H\) are mutually exclusive.
\end{itemize}
\item Venn Diagrams
\label{sec:org780a2ae}

\begin{itemize}
\item Useful for visualizing relationships among events.
\item Helpful for intuition and counterexamples.
\item Not considered rigorous proofs.
\end{itemize}
\item Event Algebra Laws
\label{sec:org9fc10ae}

\begin{enumerate}
\item Basic Identities
\label{sec:org828ea91}

\[
(E^c)^c=E
\]

\[
E\cup E^c=S
\]

\[
E\cap S=E
\]

\[
E\cap E^c=\varnothing
\]
\item Commutative Laws
\label{sec:org09c4d80}

\[
E\cup F=F\cup E
\]

\[
E\cap F=F\cap E
\]
\item Associative Laws
\label{sec:org267caae}

\[
E\cup(F\cup G)
=
(E\cup F)\cup G
\]

\[
E\cap(F\cap G)
=
(E\cap F)\cap G
\]
\item Distributive Laws
\label{sec:org2219c9a}

\[
(E\cap F)\cup H
=
(E\cup H)\cap(F\cup H)
\]

\[
(E\cup F)\cap H
=
(E\cap H)\cup(F\cap H)
\]
\item De Morgan's First Law
\label{sec:orgbb105d3}

\[
(E\cup F)^c
=
E^c\cap F^c
\]
\item De Morgan's Second Law
\label{sec:orgfe88699}

\[
(E\cap F)^c
=
E^c\cup F^c
\]
\item Useful Identity
\label{sec:org813bd48}

\[
E
=
(E\cap F)
\cup
(E\cap F^c)
\]
\end{enumerate}
\item Proof Technique: Elementwise Method
\label{sec:orgf321577}

To prove two events are equal:

\begin{enumerate}
\item Show every point of the left side belongs to the right side.
\item Show every point of the right side belongs to the left side.
\end{enumerate}

Then the two sets are equal.

Example:

\[
(E\cup F)^c
=
E^c\cap F^c
\]
\item Main Takeaways
\label{sec:org0a29767}

\begin{itemize}
\item Probability theory is built on sample spaces and events.
\item Events are simply sets of outcomes.
\item Set operations describe relationships between events.
\item Venn diagrams aid intuition.
\item Event algebra identities and De Morgan's laws are fundamental tools used throughout probability.
\end{itemize}
\end{enumerate}
\subsubsection{Additional Concepts Introduced in Exercises of Section 1.2}
\label{sec:org3addcf7}

\begin{enumerate}
\item Finite vs Infinite Sample Spaces
\label{sec:org185f631}

A sample space may contain

\begin{itemize}
\item Finite outcomes.
\item Countably infinite outcomes.
\item Uncountably infinite outcomes.
\end{itemize}

Examples:

\[
S=
\{H,TH,TTH,TTTH,\ldots\}
\]

\[
S=
\{x:-1<x<1\}
\]
\item Ordered Outcomes
\label{sec:org7774eb9}

Sometimes order matters.

Examples:

\[
HHT \neq HTH \neq THH
\]

\begin{itemize}
\item Arranging books on a shelf.
\end{itemize}
\item Cartesian Product Sample Spaces
\label{sec:orgf7e2eaf}

Experiments with several measurements are often represented as tuples.

Example:

\[
S=
\{(i,j):1\le i,j\le 6\}
\]
\item System Reliability Representation
\label{sec:org5c2e11e}

Let

\[
A_i
=
\text{component }i\text{ works}
\]

Then

\begin{itemize}
\item System works = logical combination of \(A_i\)'s.
\item System failure = complement of the working event.
\end{itemize}
\item Exactly One / Exactly Two / At Least One
\label{sec:orgaabe675}

At least one:

\[
A\cup B\cup C
\]

None:

\[
A^c\cap B^c\cap C^c
\]

Exactly one:

\[
(A\cap B^c\cap C^c)
\cup
(A^c\cap B\cap C^c)
\cup
(A^c\cap B^c\cap C)
\]

Exactly two:

\[
(A\cap B\cap C^c)
\cup
(A\cap B^c\cap C)
\cup
(A^c\cap B\cap C)
\]
\item Partition of an Event
\label{sec:orgec0d26d}

\[
E
=
(E\cap F)
\cup
(E\cap F^c)
\]

where the two pieces are mutually exclusive.
\item Infinite Unions and Intersections
\label{sec:org8662690}

For a sequence \(E_1,E_2,\ldots\),

\[
\bigcup_{i=1}^{\infty} E_i
\]

means at least one \(E_i\) occurs.

\[
\bigcap_{i=1}^{\infty} E_i
\]

means all \(E_i\) occur.

Example:

\[
E_i
=
\left(
-\frac{1}{2^i},
\frac{1}{2^i}
\right)
\]

Then

\[
\bigcup_{i=1}^{\infty} E_i
=
\left(
-\frac12,
\frac12
\right)
\]

and

\[
\bigcap_{i=1}^{\infty} E_i
=
\{0\}
\]
\item Event Sequences
\label{sec:orge878fc6}

Many exercises define

\[
A_1,A_2,A_3,\ldots
\]

and ask about

\begin{itemize}
\item infinitely many \(A_i\) occur,
\item eventually an \(A_i\) occurs,
\item all \(A_i\) occur.
\end{itemize}
\item Mutually Exclusive vs Complementary
\label{sec:org29192d6}

Mutually Exclusive:

\[
A\cap B=\varnothing
\]

Complementary:

\[
A\cap B=\varnothing
\]

and

\[
A\cup B=S
\]

Every complementary pair is mutually exclusive, but not every mutually exclusive pair is complementary.
\item Sample Space Design
\label{sec:org829a5e3}

A major skill in probability is choosing a useful sample space.

Example:

Detailed representation:

\[
\{BBB,BBG,BGB,GBB,BGG,GBG,GGB,GGG\}
\]

Compressed representation:

\[
\{0\text{ girls},1\text{ girl},2\text{ girls},3\text{ girls}\}
\]

Different representations may lead to different probabilities for sample points.
\item Event Translation Skill
\label{sec:orgfce92c5}

English statement → Set expression

At least one occurs:

\[
A\cup B\cup C
\]

None occurs:

\[
A^c\cap B^c\cap C^c
\]

\(A\) occurs but \(B\) does not:

\[
A\cap B^c
\]

\(A\) or \(B\) but not both:

\[
(A\cap B^c)
\cup
(A^c\cap B)
\]

Exactly one occurs:

\[
(A\cap B^c)
\cup
(A^c\cap B)
\]

Both occur:

\[
A\cap B
\]
\end{enumerate}
\end{document}
